% Generated draft: requires mathematical and source review.
\documentclass[11pt]{article}
\usepackage[margin=25mm]{geometry}
\usepackage{fontspec}
\setmainfont{Latin Modern Roman}
\usepackage{amsmath,amssymb,mathtools,unicode-math}
\setmathfont{Latin Modern Math}
\usepackage{enumitem,needspace,xurl,hyperref}
\setlength{\parindent}{0pt}
\setlength{\parskip}{5pt}
\setlength{\emergencystretch}{3em}
\newcommand{\N}{\mathbb{N}}
\newcommand{\Z}{\mathbb{Z}}
\newcommand{\Q}{\mathbb{Q}}
\newcommand{\R}{\mathbb{R}}
\newcommand{\C}{\mathbb{C}}
\newcommand{\F}{\mathbb{F}}
\newcommand{\A}{\mathbb{A}}
\newcommand{\PP}{\mathbb{P}}
\newcommand{\E}{\mathbb{E}}
\newcommand{\eps}{\varepsilon}
\newcommand{\dd}{\,\mathrm{d}}
\begin{document}
\section*{UnsolvedMath ID: 30005557}
\textbf{Source problem code:} catalogue entry
\subsection{Definitions and mathematical statement}
Let \(\overline{\mathcal M}_{1,1}\) denote the moduli stack of stable one-pointed curves of arithmetic genus one, and let \(\pi:\mathcal E\to\overline{\mathcal M}_{1,1}\) be its universal curve. The compactified base includes nodal curves, so this family is not everywhere a family of smooth elliptic curves. For nonnegative integers \(g,n,d\), the supplied statement uses \(\overline{\mathcal M}_{g,n}(\pi,d)\) for a moduli space of genus-\(g\), \(n\)-pointed stable maps of degree \(d\) into fibres of \(\pi\), with virtual fundamental class \([\overline{\mathcal M}_{g,n}(\pi,d)]^{\mathrm{vir}}\). The precise compactification, degree convention on singular fibres, and virtual-class construction are not specified in the supplied record.

The requested moving-target assertion is that
\[
\sum_{d\geq 0}q^d[\overline{\mathcal M}_{g,n}(\pi,d)]^{\mathrm{vir}}
\]
is quasi-modular as a cohomology-valued generating series. As written, its coefficients belong to different moduli spaces. A common coefficient space and a map taking each virtual class into that space are therefore necessary to formulate the assertion precisely.

The original discussion refers to classes on the moduli space of domain curves. A possible precise formulation, conditional on the missing choices, is the following. Suppose \(2g-2+n>0\), use rational cohomology, and suppose the chosen stable-map compactification carries the required virtual class and a proper morphism
\[
\operatorname{st}_d:\overline{\mathcal M}_{g,n}(\pi,d)\longrightarrow\overline{\mathcal M}_{g,n}
\]
that forgets the map and stabilizes the marked domain. Here \(\overline{\mathcal M}_{g,n}\) denotes the moduli stack of stable \(n\)-pointed curves of arithmetic genus \(g\). Define
\[
A_{g,n,d}=\operatorname{cl}\!\left((\operatorname{st}_d)_*[\overline{\mathcal M}_{g,n}(\pi,d)]^{\mathrm{vir}}\right)
\in H^*(\overline{\mathcal M}_{g,n};\Q),
\]
where \(\operatorname{cl}\) is the cycle-class map. If quasi-modularity means holomorphic quasi-modularity for \(\mathrm{SL}_2(\Z)\), the proposed assertion becomes
\[
\sum_{d\geq0}A_{g,n,d}q^d
\in H^*(\overline{\mathcal M}_{g,n};\Q)\otimes_{\Q}\Q[E_2,E_4,E_6],
\]
interpreted by Fourier expansion at \(q=0\), where \(E_2,E_4,E_6\) are the normalized Eisenstein series of weights \(2,4,6\). Equivalently, the series is a finite polynomial in these Eisenstein series with cohomology-class coefficients. This is a conditional interpretation, not a uniquely determined reconstruction: the source record does not specify the coefficient field, modular group, allowed poles, weight or depth conditions, admissible \((g,n)\), or the degree-zero convention.

The original discussion distinguishes eight variants according to fixed or moving elliptic target, stable maps or admissible covers, and cohomology or Chow classes. The supplied displayed statement selects moving targets, stable maps, and cohomology; it does not assert all eight variants. The supplied open label and verification label are registry assessments, not mathematical conclusions established by this draft.
\subsection{Short English statement}
The problem asks whether classes obtained from stable maps to a varying elliptic curve form a quasi-modular generating series when summed by degree. The intended classes apparently lie on the moduli space of domain curves. A precise statement still requires the pushforward construction, the treatment of nodal target fibres and degree zero, the range of genera and markings, and the meaning of quasi-modularity. The displayed polynomial formulation is conditional on explicit conventional choices.
\subsection{Sources}
\begin{itemize}
\item Local UnsolvedMath registry, \texttt{data/unsolvedmath/problems.json}, supplied record with permanent ID \(30005557\), title \emph{Quasi-Modularity of Elliptic-Curve Stable-Map Classes}, exact source problem code \texttt{catalogue entry}. The record lists proposed year 2023 and no named proposer.
\item Supplied report citation: \emph{Recent Trends in Algebraic Geometry} (2023), Oberwolfach Reports, DOI 10.4171/owr/2023/27, pp. 1485--1542. \url{https://doi.org/10.4171/owr/2023/27}. The supplied background describes these reports as documenting mathematical workshops and their open-problem discussions.
\item Supplied literature-source attribution: Pandharipande, \emph{Maps to a moving elliptic curve}. \url{https://doi.org/10.4171/owr/2023/27}.
\item Supplied literature source: \emph{Oberwolfach Report 27/2023}. \url{https://doi.org/10.4171/owr/2023/27}.
\item Literature assessment dated exactly \texttt{2026-08-22}: the supplied curation labels the question open and describes it as the moving-target analogue of a fixed-elliptic-target quasi-modularity theorem. It reports only genus-two and genus-three calculations as evidence. The original excerpt attributes these calculations to Carl Lian and cautiously says they suggest an affirmative answer may be possible. This is supplied dated curation, not independently verified literature assessment.
\item The registry marks the statement \texttt{corrected\_verified} and says it was reconstructed and checked against the official Oberwolfach report. That verification is supplied metadata; this draft has not independently inspected the report or verified the current status or cited claims.
\end{itemize}
\end{document}
